% 塔式法则的粗细两层平均。教学数据：Omega 有 8 个等概率样本点，
% X=(1,3,2,6,5,9,4,8)；细划分为相邻两点一组，组平均 2,4,7,6；
% 粗划分为前四点与后四点，组平均 3,6.5；总平均 4.75。
\begin{tikzpicture}[line width=.6pt,font=\fontsize{8}{10}\selectfont,>=Stealth,
  cell/.style={draw=FigureStroke,minimum height=6mm,inner sep=0pt},
  raw/.style={cell,minimum width=9mm,fill=white},
  fine/.style={cell,minimum width=18mm,fill=FigureBlueFill},
  coarse/.style={cell,minimum width=36mm,fill=FigureGreenFill},
  total/.style={cell,minimum width=72mm,fill=FigureYellowFill},
  lab/.style={FigureInk,anchor=east},
  note/.style={FigureMuted,font=\fontsize{7.5}{9}\selectfont,anchor=west}]
\foreach \i/\v in {0/1,1/3,2/2,3/6,4/5,5/9,6/4,7/8}{
  \node[raw] at (9mm*\i+4.5mm,0) {$\v$};}
\foreach \i/\v in {0/2,1/4,2/7,3/6}{
  \node[fine] at (18mm*\i+9mm,-10mm) {$\v$};}
\foreach \i/\v in {0/3,1/6.5}{
  \node[coarse] at (36mm*\i+18mm,-20mm) {$\v$};}
\node[total] at (36mm,-30mm) {$4.75$};
\node[lab] at (-2mm,0) {$X$};
\node[lab] at (-2mm,-10mm) {$\mathbb E[X\mid\mathcal G]$};
\node[lab] at (-2mm,-20mm) {$\mathbb E[X\mid\mathcal H]$};
\node[lab] at (-2mm,-30mm) {$\mathbb E X$};
\foreach \x in {9,27,45,63}{\draw[->,FigureBlue] (\x mm,-3.4mm) -- (\x mm,-6.6mm);}
\foreach \x in {18,54}{\draw[->,FigureGreen] (\x mm,-13.4mm) -- (\x mm,-16.6mm);}
\draw[->,FigureYellow] (36mm,-23.4mm) -- (36mm,-26.6mm);
\node[note] at (74mm,0) {完全信息};
\node[note] at (74mm,-10mm) {细划分：两点一组};
\node[note] at (74mm,-20mm) {粗划分：四点一组};
\node[note] at (74mm,-30mm) {平凡信息};
\draw[FigureRed,dashed,rounded corners=1.8pt] (0,-17mm) rectangle (36mm,-23mm);
\node[FigureRed,font=\fontsize{7.5}{9}\selectfont] at (18mm,-36mm) {$(2+4)/2=3=(1+3+2+6)/4$};
\end{tikzpicture}
